%! Describing the Distribution of a Quantitative Variable — Unit 1, Lesson 1.4 — Guided Notes & Practice
% THE in-class document, in the MAIN IDEAS / NOTES shape (see LESSON_SHAPE.md §1): the
% vocabulary box, then ONE two-column guidednotes table. Each row is one idea — a label on the
% left; on the right a terse definition the student completes, then a grid of short numbered
% problems with reserved work space. Problems are numbered 1..N through the whole component.
%   I DO   : the teacher fills the definition lines and works the FIRST problem of each grid.
%   WE DO  : the class works the rest of each grid, students holding the pen; the last row,
%            Guided Practice, is worked entirely together in a fresh context.
%   YOU DO : nothing on this page — ../homework/ IS the individual practice, started in the
%            last minutes of the period. No independent set, no reflection box, no objectives
%            box (the targets are on the cover).
% Vocabulary is GIVEN here, not discovered. The crux is THE TAIL row — the name of a skew
% follows the TAIL, not the pile (EK UNC-1.H.3); problem 15 is the crux and problem 21 is the
% crux transferred. Problem 6 is the trap: one number is never a description.
%
% THE DATA
%   Line A (20 waits, min): 3x4  4x8  5x5  6x3   -> median 4, every wait between 3 and 6
%   Line B (20 waits, min): 1x4  2x3  3x2  4x2  5x3  6x2  7x2  11  13
%                                                -> median 4, waits from 1 to 13, gap 7-11
%   Identical center, completely different experience. That is what opens the lesson.
%   Picture I   home prices, 40 homes, $100k bins:  6 14 10 5 3 1 1   (long right tail)
%   Picture II  quiz scores, 32 students, 10-pt bins: 1 2 5 11 13     (long left tail)
%   Picture III club heights, 40 adults, 2-in bins: 2 5 9 4 0 0 4 9 5 2
%                       -> symmetric AND bimodal, empty stretch 66-70 in. A single "typical
%                          height" of about 68 in belongs to nobody in the club.
%   Guided Practice: 30 students' hours of sleep: 4 | 6 6.5x2 7x4 7.5x5 8x7 8.5x6 9x3 9.5
%                       -> skewed left, peak and median at 8 h, outlier at 4 h, gap 4-6 h.
%
% PAGE LOCKSTEP with notes_key/: every \blank{} here becomes a short \ans{} there, every
% \termblank becomes a \termans, and every \pcell{n}{...}{H}{} gets its answer in the fourth
% argument. Heights and blank widths are sized to the answers so the two files set identically.
% A guidednotes row cannot break across pages — a long idea is split into sub-rows (a `\\` with
% no \hline, then `& ...`) so the table can break between them.
% NO SKETCH QUESTIONS: every display is pre-drawn in TikZ and read.
\documentclass[12pt]{article}
\usepackage{apstats-article}
\usepackage{apstats-boxes}

\begin{document}

\pageheader{Unit 1, Lesson 1.4}{Guided Notes \& Practice}
% NAMESTRIP: no \namedateperiod here — the name row lives on the cover only.

\vspace{2pt}

\boxguard[16]
\begin{vocabbox}
\small
Fill in each term \textbf{as we name it} in the notes below.
\par
\termblank{Unimodal, bimodal, uniform}
\par
\termblank{Symmetric}
\par
\termblank{Outlier}
\par
\termblank{Gap and cluster}
\par
\termblank{Skewed right, skewed left}
\par
\end{vocabbox}

\small
\begin{guidednotes}
% ── ROW 1: center and variability — the two checkout lines ───────────────────
\mainidea[Center \&]{Variability} &
A supermarket timed 20 customers in each of its two checkout lines, drawn on the same scale.
\par\vspace{2pt}\noindent\hfil
\begin{tikzpicture}[scale=0.62]
  % ── Line A: tightly bunched ──
  \draw[->, thick] (-0.2,0) -- (6.0,0);
  \foreach \x/\lab in {0/0, 0.7/2, 1.4/4, 2.1/6, 2.8/8, 3.5/10, 4.2/12, 4.9/14, 5.6/16} {
    \draw (\x,-0.07) -- (\x,0.07); \node[below, font=\tiny] at (\x,-0.1) {\lab};}
  \foreach \x/\n in {1.05/4, 1.4/8, 1.75/5, 2.1/3} {
    \foreach \k in {1,...,\n} { \fill[navy] (\x,{0.2*\k}) circle (0.07); }}
  \node[font=\tiny] at (2.8,-0.85) {Minutes waited};
  \node[font=\bfseries\small] at (2.8,2.05) {Line A};
  % ── Line B: widely spread, two customers stranded ──
  \begin{scope}[xshift=7.0cm]
    \draw[->, thick] (-0.2,0) -- (6.0,0);
    \foreach \x/\lab in {0/0, 0.7/2, 1.4/4, 2.1/6, 2.8/8, 3.5/10, 4.2/12, 4.9/14, 5.6/16} {
      \draw (\x,-0.07) -- (\x,0.07); \node[below, font=\tiny] at (\x,-0.1) {\lab};}
    \foreach \x/\n in {0.35/4, 0.7/3, 1.05/2, 1.4/2, 1.75/3, 2.1/2, 2.45/2, 3.85/1, 4.55/1} {
      \foreach \k in {1,...,\n} { \fill[goldacc] (\x,{0.2*\k}) circle (0.07); }}
    \node[font=\tiny] at (2.8,-0.85) {Minutes waited};
    \node[font=\bfseries\small] at (2.8,2.05) {Line B};
  \end{scope}
\end{tikzpicture}\hfil\par\vspace{3pt}
The \textbf{center} of a distribution is its middle, typical value --- where the data
\blank{1.2cm} up. The \textbf{variability}, or \textbf{spread}, is how far the values reach
\blank{2.6cm}. Next lesson puts numbers on both; today we describe them by eye.
\notesprompt{Read the two lines --- count in from both ends.}
\begin{probgrid}{|Y|Y|Y|}
\pcell{1}{Middle wait, Line A}{0.9cm}{} &
\pcell{2}{Middle wait, Line B}{0.9cm}{} &
\pcell{3}{Line A: lowest to highest wait}{0.9cm}{} \\ \hline
\pcell{4}{Line B: lowest to highest wait}{1.0cm}{} &
\pcell{5}{Which line has more spread?}{1.0cm}{} &
\pcell{6}{``The typical wait is 4 minutes in both lines.'' True? Useful?}{1.0cm}{} \\ \hline
\end{probgrid}
\\ \hline
% ── ROW 2: shape — count the peaks, then check the mirror ────────────────────
\mainidea[Shape]{Peaks \& Mirror} &
Three histograms: \textbf{I} sale prices of the 40 homes sold in one town last year;
\textbf{II} scores of 32 students on an easy quiz; \textbf{III} heights of the 40 adults in a
mixed volleyball club.
\par\vspace{2pt}\noindent\hfil
\begin{tikzpicture}[scale=0.58]
  % ── I: home prices, long right tail ──
  \foreach \i/\c in {0/6, 1/14, 2/10, 3/5, 4/3, 5/1, 6/1} {
    \fill[navy] ({\i*0.5},0) rectangle ({(\i+1)*0.5},{\c*0.14});
    \draw[white, very thin] ({\i*0.5},0) rectangle ({(\i+1)*0.5},{\c*0.14});}
  \draw[->, thick] (0,0) -- (0,2.5);
  \draw[->, thick] (0,0) -- (3.9,0);
  \foreach \y/\lab in {0/0, 0.7/5, 1.4/10} {
    \draw (-0.07,\y) -- (0.07,\y); \node[left, font=\tiny, xshift=-1pt] at (-0.07,\y) {\lab};}
  \foreach \x/\lab in {0/100, 1/300, 2/500, 3/700} {
    \draw (\x,-0.07) -- (\x,0.07); \node[below, font=\tiny] at (\x,-0.1) {\lab};}
  \node[font=\tiny] at (1.75,-0.85) {Sale price (\$1000s)};
  \node[font=\bfseries\small] at (1.75,2.65) {I};
  % ── II: quiz scores, long left tail ──
  \begin{scope}[xshift=4.9cm]
    \foreach \i/\c in {0/1, 1/2, 2/5, 3/11, 4/13} {
      \fill[goldacc] ({\i*0.7},0) rectangle ({(\i+1)*0.7},{\c*0.14});
      \draw[white, very thin] ({\i*0.7},0) rectangle ({(\i+1)*0.7},{\c*0.14});}
    \draw[->, thick] (0,0) -- (0,2.5);
    \draw[->, thick] (0,0) -- (3.9,0);
    \foreach \y/\lab in {0/0, 0.7/5, 1.4/10} {
      \draw (-0.07,\y) -- (0.07,\y); \node[left, font=\tiny, xshift=-1pt] at (-0.07,\y) {\lab};}
    \foreach \x/\lab in {0/50, 1.4/70, 2.8/90} {
      \draw (\x,-0.07) -- (\x,0.07); \node[below, font=\tiny] at (\x,-0.1) {\lab};}
    \node[font=\tiny] at (1.75,-0.85) {Quiz score};
    \node[font=\bfseries\small] at (1.75,2.65) {II};
  \end{scope}
  % ── III: club heights, symmetric with two humps and an empty stretch ──
  \begin{scope}[xshift=9.8cm]
    \foreach \i/\c in {0/2, 1/5, 2/9, 3/4, 4/0, 5/0, 6/4, 7/9, 8/5, 9/2} {
      \fill[navy] ({\i*0.35},0) rectangle ({(\i+1)*0.35},{\c*0.14});
      \draw[white, very thin] ({\i*0.35},0) rectangle ({(\i+1)*0.35},{\c*0.14});}
    \draw[->, thick] (0,0) -- (0,2.5);
    \draw[->, thick] (0,0) -- (3.9,0);
    \foreach \y/\lab in {0/0, 0.7/5, 1.4/10} {
      \draw (-0.07,\y) -- (0.07,\y); \node[left, font=\tiny, xshift=-1pt] at (-0.07,\y) {\lab};}
    \foreach \x/\lab in {0/58, 1.4/66, 2.8/74} {
      \draw (\x,-0.07) -- (\x,0.07); \node[below, font=\tiny] at (\x,-0.1) {\lab};}
    \node[font=\tiny] at (1.75,-0.85) {Height (inches)};
    \node[font=\bfseries\small] at (1.75,2.65) {III};
  \end{scope}
\end{tikzpicture}\hfil\par\vspace{3pt}
\textbf{Count the peaks first.} One main peak: \blank{2.0cm}. Two clear peaks: \blank{1.9cm}.
All bars about equal, no peak: \blank{1.7cm}. \textbf{Then check the mirror.} A distribution
is \textbf{symmetric} when its left half is roughly a \blank{2.4cm} of the right. Pictures I
and II instead \emph{lean}: a pile at one end, a few stragglers trailing toward the other.

\vspace{3pt}
\textbf{7.} Fill the table. The name for the lean comes in \emph{The Tail} row --- and it is
not the name most people guess.

\vspace{2pt}
\renewcommand{\arraystretch}{1.15}
\begin{tabularx}{\linewidth}{@{} c X X @{}}
\rowcolor{sky}
\textbf{Picture} & \textbf{Most of the values sit\ldots} &
\textbf{The few stragglers run\ldots} \\
\hline
\textbf{I} & \blank{3.0cm} & \blank{3.0cm} \\
\textbf{II} & \blank{3.0cm} & \blank{3.0cm} \\
\textbf{III} & \blank{3.0cm} & \blank{3.0cm} \\
\end{tabularx}
\notesprompt{Name what you see.}
\begin{probgrid}{|Y|Y|Y|}
\pcell{8}{Peaks in Picture I --- the word?}{0.9cm}{} &
\pcell{9}{Peaks in Picture III --- the word?}{0.9cm}{} &
\pcell{10}{The symmetric picture}{0.9cm}{} \\ \hline
\end{probgrid}
\\ \hline
% ── ROW 3: unusual features and the four-part description ────────────────────
\mainidea[Unusual]{Features} &
An \textbf{outlier} is a value that sits unusually far \blank{1.3cm} the rest of the data.
A \textbf{gap} is a stretch of the number line where \blank{1.6cm} was observed. A
\textbf{cluster} is a group of values bunched together and set off from the rest by a
\blank{1.0cm}.

\vspace{4pt}
\fcolorbox{navy}{sky}{\parbox{\dimexpr\linewidth-2\fboxsep-2\fboxrule\relax}{\centering
\textbf{Four things or nothing --- \blank{1.5cm}:} \textbf{S}hape $\cdot$ \textbf{O}utliers and
other unusual features $\cdot$ \textbf{C}enter $\cdot$ \textbf{S}pread\par
--- every one of the four stated \blank{2.0cm}, naming the variable and its units.}}

\vspace{3pt}
\textbf{In context} means a stranger could use the sentence.
\notesprompt{Back to Line B and Picture III.}
\begin{probgrid}{|Y|Y|}
\pcell{11}{Line B's outliers sit at \ldots; its gap runs from \ldots\ to \ldots}{1.0cm}{} &
\pcell{12}{Clusters in Picture III --- how many, and where?}{1.0cm}{} \\ \hline
\pcell{13}{``Center about 4, spread 1 to 13.'' Rewrite it so a stranger could use it.}{1.5cm}{} &
\pcell{14}{These 20 customers were timed on one day. What does that say about the store's lines
in general?}{1.5cm}{} \\ \hline
\end{probgrid}
\\ \hline
% ── ROW 4: the tail decides — THE CRUX ───────────────────────────────────────
\mainidea[Skewed which way?]{The Tail} &
Count Picture I's homes. Below \$400{,}000: $6 + 14 + 10 = $ \blank{0.9cm} of the 40. Above
\$600{,}000: \blank{0.9cm} of the 40 --- yet the picture stretches past \$700{,}000, because
two houses really did sell for that. The pile is on the \blank{1.0cm}; the long thin tail runs
to the \blank{1.2cm}. So Picture I is \textbf{skewed right}: piled \emph{left}, named \emph{right}.

\vspace{4pt}
\fcolorbox{navy}{sky}{\parbox{\dimexpr\linewidth-2\fboxsep-2\fboxrule\relax}{\centering
\textbf{The name of a skew follows the \blank{1.1cm}, not the \blank{1.1cm}.}\par
Skewed \textbf{right}: the longer tail runs right, so most values are \blank{1.4cm} and a few
are much larger. Skewed \textbf{left}: the longer tail runs \blank{1.1cm}.}}
\\
& \textbf{All four for Picture I:} sale prices are \blank{2.1cm} and unimodal; no gaps, and
\blank{0.9cm} homes above \$600{,}000; the middle home sold for about \$\blank{1.6cm}; prices
run from about \$100{,}000 to \$\blank{1.6cm}.
\notesprompt{Point at the tail, then say the name.}
\begin{probgrid}{|Y|Y|Y|}
\pcell{15}{Picture II: which side is the pile on? Skewed which way?}{1.5cm}{} &
\pcell{16}{``Picture II is skewed right --- its tallest bars are on the right.'' What got read
instead of the tail?}{1.5cm}{} &
\pcell{17}{Picture III is symmetric. Is that the whole shape? Where is its gap?}{1.5cm}{} \\ \hline
\pcell{18}{``The typical height in this club is about 68 inches.'' True of how many
adults?}{1.5cm}{} &
\pcell{19}{Line B's 13-minute customer: delete her, or check her? Why?}{1.5cm}{} &
\pcell{20}{Name a variable you expect to be skewed right, and say why.}{1.5cm}{} \\ \hline
\end{probgrid}
\\ \hline
% ── ROW 5: guided practice — WE DO, together, fresh context ──────────────────
\mainidea[Guided practice]{The Sleep Survey} &
Thirty students reported how many hours they slept on a school night.
\par\vspace{2pt}\noindent\hfil
\begin{tikzpicture}[scale=0.62]
  \draw[->, thick] (-0.2,0) -- (8.2,0);
  \foreach \x/\lab in {0.6/4, 1.8/5, 3.0/6, 4.2/7, 5.4/8, 6.6/9, 7.8/10} {
    \draw (\x,-0.07) -- (\x,0.07); \node[below, font=\tiny] at (\x,-0.1) {\lab};}
  \foreach \x/\n in {0.6/1, 3.0/1, 3.6/2, 4.2/4, 4.8/5, 5.4/7, 6.0/6, 6.6/3, 7.2/1} {
    \foreach \k in {1,...,\n} { \fill[navy] (\x,{0.2*\k}) circle (0.07); }}
  \node[font=\tiny] at (4.0,-0.85) {Hours of sleep on a school night};
\end{tikzpicture}\hfil\par
\notesprompt{We work these together.}
\begin{probgrid}{|Y|Y|}
\pcell{21}{Which way does it lean? Point at the tail, then name it.}{1.0cm}{} &
\pcell{22}{Students shown; peaks; the value that sits apart, and its gap.}{1.0cm}{} \\ \hline
\pcell{23}{What does the longer tail mean for \emph{these students}, not for the graph?}{1.6cm}{} &
\pcell{24}{All four, in context: shape, unusual features, center, spread --- variable and
units.}{2.2cm}{} \\ \hline
\end{probgrid}
\\ \hline
\end{guidednotes}

% ── THE NOTES END AT GUIDED PRACTICE ─────────────────────────────────────────
% There is deliberately no "you do" here: ../homework/ is the individual practice,
% started in class in the last 8 minutes while the teacher circulates.

\end{document}
